\par\medskip
\noindent\textbf{Teorema 6.4.}
Sea $X_0$ un $k$-esquema de tipo finito cuyos puntos no lisos sean
aislados, y sea $(R,\mathfrak X)$ una deformación versal formal de
$X_0$. Denotemos por $\{z_1,\ldots,z_m\}$ el conjunto de puntos no
lisos de $X_0$. Supongamos que
\[
H^2(X_0,\underline D^0_{X_0})=0
\]
y que todo $z_i$ es una singularidad cuadrática no degenerada
racional sobre $k$. Entonces:

\smallskip
\noindent\textbf{(1)} $R$ es formalmente liso sobre $\Lambda$; es
decir, $R$ es un anillo de series formales sobre $\Lambda$.

\smallskip
\noindent\textbf{(2)} Hay una descomposición en unión disjunta
\[
J_{\mathfrak X}(\mathfrak X|R)
=\coprod_{1\leq i\leq m}
 J_{\mathfrak X_{(z_i)}}(\mathfrak X_{(z_i)}|R),
\]
y el morfismo canónico
\[
J_{\mathfrak X_{(z_i)}}(\mathfrak X_{(z_i)}|R)
 \xrightarrow{\ \sim\ }
J_R(\mathfrak X_{(z_i)}|R)
\]
es un isomorfismo.

\smallskip
\noindent\textbf{(3)} Cada ideal
$I_R(\mathfrak X_{(z_i)}|R)$ es principal, al igual que
\[
I_R(\mathfrak X|R)
=\prod_{1\leq i\leq m}I_R(\mathfrak X_{(z_i)}|R).
\]
Si $I_R(\mathfrak X_{(z_i)}|R)=(t_i)$, entonces
$t_1,\ldots,t_m$ pueden completarse hasta un sistema de generadores
formales libres de $R$ sobre $\Lambda$.

\smallskip
\noindent\textbf{(4)} El morfismo
$\mathfrak X\longrightarrow\operatorname{Spec}(\Lambda)$ es
formalmente liso; equivalentemente,
\[
J_{\mathfrak X}(\mathfrak X|\Lambda)=\varnothing.
\]

%%%%%%%%%%  scan idx 139  |  folio 128  %%%%%%%%%%

\smallskip
\noindent\textbf{(5)} Si $X_0$ es una curva completa irreducible,
entonces
\[
\dim\underline M_{X_0}=3g-3+a,
\]
donde
\[
g=\dim_kH^1(X_0,\mathcal O_{X_0}),
\qquad
a=\dim_kH^0(X_0,\underline D^0_{X_0}).
\]

\smallskip
\noindent\textit{Demostración.}
La afirmación (1) se deduce de (2) y (3). La primera afirmación de (2)
se deduce de 5.7(2). Para cada $i$, sea $(R_i,\mathfrak Y_i)$ una
deformación versal formal de
$\operatorname{Spec}(\widehat{\mathcal O}_{X_0,z_i})$. Como
$(R,\mathfrak X_{(\widehat z_i)})$ es una deformación de este tipo,
existe un homomorfismo local continuo de $\Lambda$-álgebras
\[
\varphi_i:R_i\longrightarrow R
\]
y un diagrama cartesiano
\[
\begin{tikzcd}[column sep=huge,row sep=large]
\mathfrak X_{(\widehat z_i)}
  \arrow[r] \arrow[d]
& \mathfrak Y_i \arrow[d]\\
\operatorname{Spf}(R)
  \arrow[r,"{\operatorname{Spf}(\varphi_i)}"']
& \operatorname{Spf}(R_i).
\end{tikzcd}
\]
Por consiguiente,
\[
\begin{aligned}
J_{\mathfrak X_{(z_i)}}(\mathfrak X_{(z_i)}|R)
&=J_{\mathfrak X_{(\widehat z_i)}}
   (\mathfrak X_{(\widehat z_i)}|R)\\
&\simeq
J_{\mathfrak Y_i}(\mathfrak Y_i|R_i)
 \times_{\operatorname{Spf}(R_i)}\operatorname{Spf}(R),
\end{aligned}
\]
y
\[
\begin{aligned}
J_R(\mathfrak X_{(z_i)}|R)
&=J_R(\mathfrak X_{(\widehat z_i)}|R)\\
&\simeq
J_{R_i}(\mathfrak Y_i|R_i)
 \times_{\operatorname{Spf}(R_i)}\operatorname{Spf}(R).
\end{aligned}
\]
Por el Lema 6.2(1),
\[
J_{\mathfrak Y_i}(\mathfrak Y_i|R_i)
 \longrightarrow J_{R_i}(\mathfrak Y_i|R_i)
\]
es un isomorfismo, lo cual demuestra la afirmación restante de (2).

Por el Lema 6.2(2),
\[
I_{R_i}(\mathfrak Y_i|R_i)=(u_i),
\qquad R_i=\Lambda[[u_i]].
\]
Se sigue que
\[
I_R(\mathfrak X_{(z_i)}|R)
=I_R(\mathfrak X_{(\widehat z_i)}|R)
=(t_i),
\qquad t_i=\varphi_i(u_i),
\]
y por tanto
\[
I_R(\mathfrak X|R)
=\prod_{1\leq i\leq m}I_R(\mathfrak X_{(z_i)}|R)
=\left(\prod_{i=1}^m t_i\right).
\]
Por 4.13, el morfismo natural
\[
\underline M_{X_0}\longrightarrow
\underline M_{X_0,(\widehat z_1)}\times\cdots\times
\underline M_{X_0,(\widehat z_m)}
\]
es liso. Por tanto,
\[
R_1\widehat\otimes_\Lambda\cdots
 \widehat\otimes_\Lambda R_m\longrightarrow R
\]
es formalmente liso. Así, $t_1,\ldots,t_m$ se completan hasta un
sistema de generadores formales libres de $R$ sobre $\Lambda$, lo que
demuestra (1) y (3).

%%%%%%%%%%  scan idx 140  |  folio 129  %%%%%%%%%%

Para (4), la proyección
\[
\mathfrak X_{(\widehat z_i)}
 \xrightarrow{\ \sim\ }
\mathfrak Y_i\times_{\operatorname{Spf}(R_i)}
 \operatorname{Spf}(R)
 \longrightarrow\mathfrak Y_i
\]
es formalmente lisa porque
$\operatorname{Spf}(R)\sisi{\overset{\ \operatorname{Spf}(\varphi_i)\ }{\relbar\joinrel\relbar\joinrel\relbar\joinrel\relbar}}{\longrightarrow}\operatorname{Spf}(R_i)$ es
formalmente liso. El Lema 6.2(3) muestra que
$\mathfrak Y_i\longrightarrow\operatorname{Spec}(\Lambda)$ es
formalmente liso. Por consiguiente, todo
$\mathfrak X_{(\widehat z_i)}\longrightarrow
\operatorname{Spec}(\Lambda)$ es formalmente liso. Como los $z_i$
son precisamente los puntos no lisos de $X_0$, se sigue que
$\mathfrak X\longrightarrow\operatorname{Spec}(\Lambda)$ es
formalmente liso.
